\documentclass[11pt,a4paper]{article}
\usepackage[T1]{fontenc}
\usepackage[utf8]{inputenc}
\usepackage{amsmath,amssymb,amsthm}
\usepackage[margin=1in]{geometry}
\usepackage[colorlinks=true,linkcolor=blue,urlcolor=blue]{hyperref}
\theoremstyle{plain}
\newtheorem{theorem}{Theorem}
\newtheorem{lemma}[theorem]{Lemma}
\newtheorem{proposition}[theorem]{Proposition}
\newtheorem{corollary}[theorem]{Corollary}
\theoremstyle{definition}
\newtheorem{remark}[theorem]{Remark}

\title{The empirical recurrence for OEIS A210072, proved by transfer matrix}
\author{Adrian Perez Fontelles\\ \small Independent researcher}
\date{31 August 2026}
\begin{document}
\maketitle

\begin{abstract}
OEIS A210072 counts $(n+1)\times4$ arrays over $\{0,\dots,3\}$ subject to a condition imposed on
every $2\times2$ subblock, and carries an empirical recurrence of order $71$ contributed
by R.~H.~Hardin. It is true, and it is not an empirical matter at all. The condition is local
to each subblock, so the rows of an admissible array are the vertices of a finite digraph
on $S=256$ vertices whose edges are the admissible consecutive pairs, and the count is a
number of walks: $a(n)=\tfrac1{4}\mathbf{1}^{\!\top}M^{n}\mathbf{1}$. Writing $q$
for the characteristic polynomial of the conjectured recurrence, the recurrence holds for all
$n$ exactly when $\mathbf{1}^{\!\top}M^{j}q(M)\mathbf{1}=0$ for every $j\ge0$, and the
Cayley--Hamilton theorem bounds that to $j<S$. It is therefore a finite computation in exact
integer arithmetic, and it comes out zero.
\end{abstract}

\noindent\small 2020 Mathematics Subject Classification. 05A15, 05B45, 15B36.\normalsize

\section{The sequence and the conjecture}

OEIS A210072 is ``1/4 the number of (n+1)X4 0..3 arrays with every 2X2 subblock having one or three distinct clockwise edge differences.''. It has offset $1$ and begins
\[
2908, 183628, 11333893, 707627792, 44179487805, 2759914868891, 172422974616663, 10772379159246366,\ \dots
\]
The entry carries the following comment:
\begin{quote}\small
Empirical: a(n) = 105*a(n-1) -2719*a(n-2) -15451*a(n-3) +1390943*a(n-4) -7924538*a(n-5) -266634929*a(n-6) +2783228062*a(n-7) +26184293659*a(n-8) -399963589595*a(n-9) -1390685894189*a(n-10) +33792993132607*a(n-11) +30902632289869*a(n-12) -1890381555307832*a(n-13) +822619990486604*a(n-14) +74406664819656249*a(n-15) -89835422690076184*a(n-16) -2138656368143245036*a(n-17) +3544878797500349305*a(n-18) +46008317343363322887*a(n-19) -87903800402488590784*a(n-20) -753473801980270591181*a(n-21) +1535436320348069965003*a(n-22) +9505348644817473245840*a(n-23) -19817177440215843837829*a(n-24) -93126842954944481973745*a(n-25) +194008472781630232252845*a(n-26) +712417629951787189697006*a(n-27) -1464351302320119255078001*a(n-28) -4269039597123343663401254*a(n-29) +8611866713944156678864641*a(n-30) +20062337369297929637737396*a(n-31) -39721870514776593763639623*a(n-32) -73883618013286155451079859*a(n-33) +144186462856893414613730955*a(n-34) +212594575285455273050900740*a(n-35) -412162806244228323032035765*a(n-36) -475362727557671730453508543*a(n-37) +926067583198793638831571069*a(n-38) +818943669820992562798006291*a(n-39) -1628407201463016755238157088*a(n-40) -1073377817747961549292055653*a(n-41) +2225907870727063061101510830*a(n-42) +1050655346169817900693995129*a(n-43) -2343867586300105570589047976*a(n-44) -746446528111441994989443963*a(n-45) +1879795853747926271266149480*a(n-46) +366426635285816505070213167*a(n-47) -1132874535089578641767501432*a(n-48) -111323955733240541695109281*a(n-49) +505194638908821997108469349*a(n-50) +12927626382477733740240979*a(n-51) -163887799475213923454810367*a(n-52) +4251111577638305441346089*a(n-53) +37966681387926744734306902*a(n-54) -2249309141628534316063828*a(n-55) -6153601732703873799844070*a(n-56) +476254714828361557941333*a(n-57) +681177457210232099839986*a(n-58) -56897891948249201993222*a(n-59) -49910567815506758541560*a(n-60) +4016897111990219118375*a(n-61) +2316946754025069893095*a(n-62) -163269106694177521346*a(n-63) -63879794902532341859*a(n-64) +3536992747961553187*a(n-65) +943996081685851752*a(n-66) -35450092561754998*a(n-67) -6289331583619410*a(n-68) +135657995058528*a(n-69) +14840999586360*a(n-70) -115072392784*a(n-71)
\end{quote}
As of the ``Last modified'' line on the live entry (Oct 03 2025 15:40:58, revision 9) the statement is
still recorded as empirical, and nothing on the entry records it as settled either way.

Written out, the claim is
\begin{equation}\label{eq:conj}
a(n)\;=\;105\,a(n-1) - 2719\,a(n-2) - 15451\,a(n-3) + 1390943\,a(n-4) - 7924538\,a(n-5) - 266634929\,a(n-6) + 2783228062\,a(n-7) + 26184293659\,a(n-8) - 399963589595\,a(n-9) - 1390685894189\,a(n-10) + 33792993132607\,a(n-11) + 30902632289869\,a(n-12) - 1890381555307832\,a(n-13) + 822619990486604\,a(n-14) + 74406664819656249\,a(n-15) - 89835422690076184\,a(n-16) - 2138656368143245036\,a(n-17) + 3544878797500349305\,a(n-18) + 46008317343363322887\,a(n-19) - 87903800402488590784\,a(n-20) - 753473801980270591181\,a(n-21) + 1535436320348069965003\,a(n-22) + 9505348644817473245840\,a(n-23) - 19817177440215843837829\,a(n-24) - 93126842954944481973745\,a(n-25) + 194008472781630232252845\,a(n-26) + 712417629951787189697006\,a(n-27) - 1464351302320119255078001\,a(n-28) - 4269039597123343663401254\,a(n-29) + 8611866713944156678864641\,a(n-30) + 20062337369297929637737396\,a(n-31) - 39721870514776593763639623\,a(n-32) - 73883618013286155451079859\,a(n-33) + 144186462856893414613730955\,a(n-34) + 212594575285455273050900740\,a(n-35) - 412162806244228323032035765\,a(n-36) - 475362727557671730453508543\,a(n-37) + 926067583198793638831571069\,a(n-38) + 818943669820992562798006291\,a(n-39) - 1628407201463016755238157088\,a(n-40) - 1073377817747961549292055653\,a(n-41) + 2225907870727063061101510830\,a(n-42) + 1050655346169817900693995129\,a(n-43) - 2343867586300105570589047976\,a(n-44) - 746446528111441994989443963\,a(n-45) + 1879795853747926271266149480\,a(n-46) + 366426635285816505070213167\,a(n-47) - 1132874535089578641767501432\,a(n-48) - 111323955733240541695109281\,a(n-49) + 505194638908821997108469349\,a(n-50) + 12927626382477733740240979\,a(n-51) - 163887799475213923454810367\,a(n-52) + 4251111577638305441346089\,a(n-53) + 37966681387926744734306902\,a(n-54) - 2249309141628534316063828\,a(n-55) - 6153601732703873799844070\,a(n-56) + 476254714828361557941333\,a(n-57) + 681177457210232099839986\,a(n-58) - 56897891948249201993222\,a(n-59) - 49910567815506758541560\,a(n-60) + 4016897111990219118375\,a(n-61) + 2316946754025069893095\,a(n-62) - 163269106694177521346\,a(n-63) - 63879794902532341859\,a(n-64) + 3536992747961553187\,a(n-65) + 943996081685851752\,a(n-66) - 35450092561754998\,a(n-67) - 6289331583619410\,a(n-68) + 135657995058528\,a(n-69) + 14840999586360\,a(n-70) - 115072392784\,a(n-71)
\end{equation}
for all $n>70$.

\section{The condition is local, so the rows form a digraph}

Write an array of the shape in question as its list of rows
$r^{(1)},r^{(2)},\dots$, each an element of
$V=\{0,1,\dots,3\}^{4}$, so $|V|=S=256$. A $2\times2$ subblock of the array
occupies two consecutive rows and two consecutive columns; if the two rows are $r$
and $s$ (in that order) and the two columns are numbered $j,j+1$,
the subblock is
\[
\begin{pmatrix}\alpha&\beta\\\gamma&\delta\end{pmatrix}=\begin{pmatrix}r_j&r_{j+1}\\s_j&s_{j+1}\end{pmatrix}.
\]
The entry's requirement on that subblock is
\begin{equation}\label{eq:loc}
\#\{\beta-\alpha,\ \delta-\beta,\ \gamma-\delta,\ \alpha-\gamma\}\in\{1,3\},
\end{equation}
a condition on the four entries alone.

For $r,s\in V$ declare $r\to s$ an edge of a digraph $G$ when, for every
$1\le j\le 4-1$,
\begin{itemize}
\item the four entries above satisfy \eqref{eq:loc};

\end{itemize}
Let $M\in\{0,1\}^{S\times S}$ be the adjacency matrix of $G$ and $\mathbf{1}$ the
all-ones vector.

\begin{lemma}\label{lem:walk}
The arrays counted by the entry with $L$ rows are exactly the walks
$r^{(1)}\to r^{(2)}\to\cdots\to r^{(L)}$ of length $L-1$ in $G$. Consequently
\[
a(n)\;=\;\tfrac1{4}\mathbf{1}^{\!\top}M^{n}\mathbf{1} .
\]
\end{lemma}

\begin{proof}
An array with $L$ rows and $4$ columns is precisely a sequence
$r^{(1)},\dots,r^{(L)}$ of elements of $V$; conversely any such sequence is an array.
Every $2\times2$ subblock lies in two consecutive rows $r^{(t)},r^{(t+1)}$ and two
consecutive columns, so the array satisfies the entry's condition on all of its subblocks
if and only if each consecutive pair $\bigl(r^{(t)},r^{(t+1)}\bigr)$ satisfies it for
every column pair --- which is exactly the edge relation of $G$. The number of walks of
length $L-1$, summed over all starting and ending vertices, is
$\mathbf{1}^{\!\top}M^{L-1}\mathbf{1}$, since $(M^{k})_{r,s}$ counts the walks of
length $k$ from $r$ to $s$. The entry's index $n$ corresponds to $L=n+1$ rows. The entry counts $\tfrac1{4}$ of those arrays, a constant factor that passes through every step below unchanged.
\end{proof}

\section{The criterion}

Let $q(t)=t^{71}-\sum_i c_i\,t^{71-i}$ be the characteristic polynomial of
\eqref{eq:conj}, with the $c_i$ read off that equation.

\begin{lemma}\label{lem:crit}
For every $n$ with $n+1-1\ge71$,
\[
a(n)-\sum_i c_i\,a(n-i)\;=\;\tfrac1{4}\mathbf{1}^{\!\top}M^{\,n+1-1-71}\,q(M)\,\mathbf{1} .
\]
Hence \eqref{eq:conj} holds from that point on if and only if
$\mathbf{1}^{\!\top}M^{j}q(M)\mathbf{1}=0$ for every $j\ge0$; and it is enough to check
$j=0,1,\dots,S-1$.
\end{lemma}

\begin{proof}
By Lemma~\ref{lem:walk}, $a(n-i)=\tfrac1{4}\mathbf{1}^{\!\top}M^{\,n+1-1-i}\mathbf{1}$, so
\[
a(n)-\sum_i c_i a(n-i)
=\tfrac1{4}\mathbf{1}^{\!\top}\Bigl(M^{m}-\sum_i c_i M^{m-i}\Bigr)\mathbf{1}
=\tfrac1{4}\mathbf{1}^{\!\top}M^{\,m-71}\,q(M)\,\mathbf{1}
\]
with $m=n+1-1$, which is the stated identity; the scalar $\tfrac1{4}$ never vanishes, so
it does not affect whether the left side is zero. The equivalence follows by letting
$m-71=j$ run over $\ge0$. For the last clause put $w=q(M)\mathbf{1}$. By the
Cayley--Hamilton theorem $M^{S}$ is an integer combination of $I,M,\dots,M^{S-1}$, so
$\mathbf{1}^{\!\top}M^{j}w$ for $j\ge S$ is the corresponding combination of the earlier
values; if those all vanish, so do all the rest.
\end{proof}

\section{The computation}

\begin{theorem}
The recurrence \eqref{eq:conj} holds for every $n>70$, which is the statement of
Section~1.
\end{theorem}

\begin{proof}
The digraph was constructed from \eqref{eq:loc} by a depth-first scan across the $4$
columns: the edge condition is a conjunction of one constraint per consecutive column
pair, so the successors of a row $r$ are enumerated column by column without testing all
$S^2$ pairs. The vector $w=q(M)\mathbf{1}\in\mathbb{Z}^{S}$ was then formed by $71$
matrix--vector products in exact integer arithmetic, and $\mathbf{1}^{\!\top}M^{j}w$ was
evaluated for $j=0,1,\dots$, again exactly. Every one of those integers vanishes from the
index corresponding to $n=70+1$ onwards, and the run continued until $w$ itself was the
zero vector, which settles every larger $j$ at once. By Lemma~\ref{lem:crit} the recurrence
holds for every $n>70$.
\end{proof}

\begin{remark}
The argument gives more than the recurrence: it shows the sequence is $C$-finite with
characteristic polynomial dividing that of $M$, so it has a rational generating function
whose denominator divides $\det(I-xM)$. The conjectured recurrence is one consequence; the
minimal one is a divisor of $q$.
\end{remark}

\section{Verification}

Three checks, each independent of the algebra above.

First, the digraph was built directly from the entry's stated condition and the walk counts
$\tfrac1{4}\mathbf{1}^{\!\top}M^{L-1}\mathbf{1}$ compared against the entry's DATA at
every one of the $12$ published indices; the values agree exactly. That is what ties
the digraph to the sequence: a model that reproduced only some of the published terms would
be the wrong model, and would have been rejected. In particular it is what rules out a
misreading of the entry's wording, since a different reading of the condition gives a
different digraph and different counts.

Second, the recurrence \eqref{eq:conj} was evaluated on those published terms in exact
integer arithmetic, with no matrices involved, and vanishes wherever it applies.

Third, the annihilation test of Lemma~\ref{lem:crit} was run in exact integer arithmetic
until the working vector was identically zero, not to a sampled prefix, and returned zero
throughout.

\begin{thebibliography}{9}
\bibitem{oeis} The OEIS Foundation, \emph{The On-Line Encyclopedia of Integer
Sequences}, \url{https://oeis.org}, 2026. Sequence A210072.
\bibitem{stanley} R.~P.~Stanley, \emph{Enumerative Combinatorics}, Volume 1, 2nd ed.,
Cambridge University Press, 2011. (Transfer-matrix method, Section 4.7.)
\bibitem{fs} P.~Flajolet and R.~Sedgewick, \emph{Analytic Combinatorics}, Cambridge
University Press, 2009.
\bibitem{horn} R.~A.~Horn and C.~R.~Johnson, \emph{Matrix Analysis}, 2nd ed., Cambridge
University Press, 2013.
\end{thebibliography}

\end{document}
